%!TEX root = ../main.tex
\chapter{Fundamentals}
\label{ch:fundamentals}

Scientific thinking is not a personality trait; it is a small set of
habits that anyone can practice. This chapter covers two of them:
making a claim precise, and checking it against evidence.

\section{Making a Claim Precise}
\label{sec:precise-claims}

\begin{definition}[Falsifiable claim]
A claim is \emph{falsifiable} if there exists some observation that,
were it made, would show the claim to be false.
\end{definition}

``This medicine helps some people'' is not falsifiable in any useful
sense — almost any outcome is compatible with it. ``This medicine
reduces average recovery time by at least two days, in a controlled
trial with a placebo arm'' is falsifiable: a trial could come back
showing no difference, or a difference smaller than two days.

\begin{example}
Compare ``vitamin C prevents colds'' with ``taking 1000mg of vitamin C
daily reduces the annual number of colds in adults by more than 20\%,
measured against a placebo control.'' The second version can actually
be tested — and, when it has been, the reduction found is usually much
smaller than folklore suggests.
\end{example}

\section{Checking Against Evidence}
\label{sec:evidence}

Precision alone is not enough; the claim also has to be checked
against an observation that could, in principle, have gone the other
way. A prediction confirmed only by evidence you already knew about
before making the prediction does not count — this is sometimes called
\emph{postdiction} rather than prediction, and it is far easier to get
right by accident.
